\section{Introduction}

\subsection{Purpose and Scope}

This report documents a \textbf{structured, quantified} usability-style evaluation of Hitchly, aligned with HCI reporting practice. It specifies the \textbf{goals} of testing, the \textbf{tasks} performed, the \textbf{questions} and \textbf{survey instruments} used for feedback, and the \textbf{metrics} used to summarize outcomes. Section~\ref{sec:methodology} and Section~\ref{sec:results} state clearly that aggregates are \textbf{illustrative} ($N=10$ simulated personas) rather than a live external-participant study.

The evaluation targets the mobile experience (and any supporting flows such as email verification) for McMaster-affiliated riders and drivers. It complements verification activities in the Verification and Validation (VnV) plan by focusing on learnability, efficiency, error tolerance, and satisfaction for representative users rather than on code-level correctness alone.

\subsection{Relation to Requirements}

Usability objectives align with non-functional expectations in the SRS for understandable, low-friction flows (e.g., sign-up, offering or requesting a ride, confirming a match, payment-related steps where applicable). Findings may be traced to specific FR/NFR identifiers when reporting issues or recommendations.

\subsection{Document Structure}

Section~\ref{sec:goals} states test goals. Section~\ref{sec:methodology} describes participants, setting, and procedure. Section~\ref{sec:task_scenarios} lists task scenarios. Section~\ref{sec:questions_surveys} covers moderator questions and surveys. Section~\ref{sec:metrics} defines quantitative measures and analysis. Sections~\ref{sec:results} and~\ref{sec:discussion} contain \textbf{example} results and recommendations from simulated walkthroughs. Appendix~\ref{sec:appendix} holds questionnaire text.
